%%%%%%%%%%%%%%%%%%%%%%%%%%%%%%%%%%%%%%%%%%%%%%%%%%%%%%%%%%%%%%%%%%
% Resumo em Português
%%%%%%%%%%%%%%%%%%%%%%%%%%%%%%%%%%%%%%%%%%%%%%%%%%%%%%%%%%%%%%%%%%

\chapter*{Resumo}
\addcontentsline{toc}{chapter}{Resumo}

O processamento de imagens hiperespectrais em sistemas embarcados representa 
um dos desafios computacionais mais complexos da atualidade, exigindo 
conciliação entre requisitos aparentemente conflitantes de alta performance, 
baixa latência e eficiência energética em plataformas com severas restrições 
de recursos. Esta dissertação apresenta um estudo comparativo para otimização 
computacional de processamento hiperespectral embarcado, focando na análise de 
trade-offs entre diferentes arquiteturas heterogêneas e técnicas de otimização 
através de abordagem exploratória-comparativa baseada em simulação e modelagem teórica.

A pesquisa adota abordagem metodológica que privilegia análise 
teórica e simulada sobre implementação física, justificada pela necessidade 
de investigar sistematicamente múltiplas alternativas arquiteturais sob 
condições controladas. A metodologia estrutura-se em oito componentes principais: 
catalogação sistemática de técnicas de otimização FPGA e GPU; modelagem 
teórica comparativa de performance; framework de simulação integrado; estudo 
comparativo cross-platform; protocolo de benchmarking padronizado; análise 
multi-critério baseada em metodologia AHP e análise de Pareto; design space 
exploration para identificação de configurações ótimas; e síntese de arquitetura 
heterogênea otimizada.

A revisão sistemática de 25 trabalhos publicados entre 2011-2024 revelou 
lacuna metodológica significativa: apenas 4\% dos estudos abordam sistemas 
completos integrados, enquanto 96\% focam em técnicas isoladas ou otimizações 
específicas de plataforma. Esta fragmentação impede análise sistemática de 
trade-offs e identificação de soluções ótimas para cenários específicos de 
aplicação. Como principal contribuição, desenvolve-se framework metodológico 
unificado que permite comparação objetiva entre diferentes abordagens 
arquiteturais considerando múltiplos critérios de performance.

A metodologia proposta estabelece base científica rigorosa para investigação 
de sistemas hiperespectrais embarcados, preenchendo lacuna identificada na 
literatura e proporcionando framework reprodutível para pesquisas futuras. 
Os resultados esperados incluem identificação quantitativa de trade-offs 
ssfundamentais, proposição de arquitetura heterogênea otimizada baseada em 
análise sistemática, e desenvolvimento de recomendações práticas para seleção 
de técnicas segundo requisitos específicos de aplicação em agricultura de 
precisão, monitoramento ambiental e sistemas UAV autônomos.

\textbf{Palavras-chave:} 
Processamento Hiperespectral, Arquiteturas Heterogêneas, Metodologia Comparativa, Otimização Computacional, Análise Multi-critério, Sistemas Embarcados, FPGA, GPU.
